\documentclass[12pt,a4paper]{article}

% ==================== PACKAGES ====================
\usepackage[utf8]{inputenc}
\usepackage[vietnamese]{babel}
\usepackage{amsmath,amssymb,amsthm}
\usepackage{geometry}
\usepackage{enumitem}
\usepackage[most]{tcolorbox}
\usepackage{hyperref}
\usepackage{fancyhdr}
\usepackage{graphicx}
\usepackage{tikz}
\usepackage{eso-pic}
\usepackage{xcolor}

\geometry{a4paper, margin=2cm, bottom=2.5cm}
\hypersetup{colorlinks=true, linkcolor=blue!70!black, urlcolor=blue}
\setlist[itemize]{topsep=4pt, itemsep=2pt}
\setlist[enumerate]{topsep=4pt, itemsep=2pt}

\AddToShipoutPictureFG{%
  \AtPageCenter{%
    \begin{tikzpicture}[remember picture, overlay]
      \node[rotate=45, scale=1, opacity=0.06]{\fontsize{60}{72}\selectfont\bfseries\sffamily Đặng Tấn Quí};
    \end{tikzpicture}%
  }%
}

\pagestyle{fancy}
\fancyhf{}
\fancyhead[L]{\small\textit{Giải tích 1}}
\fancyhead[R]{\small\textit{Đặng Tấn Quí}}
\fancyfoot[C]{\thepage}
\fancyfoot[L]{\footnotesize\textcopyright\ Đặng Tấn Quí}
\fancyfoot[R]{\footnotesize SĐT: 0377\,122\,210}
\renewcommand{\headrulewidth}{0.4pt}
\renewcommand{\footrulewidth}{0.4pt}

\fancypagestyle{plain}{%
  \fancyhf{}
  \fancyfoot[C]{\thepage}
  \fancyfoot[L]{\footnotesize\textcopyright\ Đặng Tấn Quí}
  \fancyfoot[R]{\footnotesize SĐT: 0377\,122\,210}
  \renewcommand{\headrulewidth}{0pt}
  \renewcommand{\footrulewidth}{0.4pt}
}

\newtcolorbox{dinhnghia}[1][]{%
  enhanced, breakable, colback=blue!3!white, colframe=blue!60!black,
  fonttitle=\bfseries, title={Định nghĩa}, #1%
}
\newtcolorbox{dinhly}[1][]{%
  enhanced, breakable, colback=green!3!white, colframe=green!50!black,
  fonttitle=\bfseries, title={Định lý}, #1%
}
\newtcolorbox{tinhchat}[1][]{%
  enhanced, breakable, colback=yellow!5!white, colframe=orange!50!black,
  fonttitle=\bfseries, title={Tính chất}, #1%
}
\newtcolorbox{phuongphap}[1][]{%
  enhanced, breakable, colback=gray!5!white, colframe=gray!50!black,
  fonttitle=\bfseries, title={Phương pháp}, #1%
}
\newtcolorbox{luuy}[1][]{%
  enhanced, breakable, colback=red!3!white, colframe=red!50!black,
  fonttitle=\bfseries, title={Lưu ý}, #1%
}
\newtcolorbox{congthuc}[1][]{%
  enhanced, breakable, colback=cyan!3!white, colframe=cyan!50!black,
  fonttitle=\bfseries, title={Công thức}, #1%
}
\newtcolorbox{vidu}[1][]{%
  enhanced, breakable, colback=violet!3!white, colframe=violet!50!black,
  fonttitle=\bfseries, title={Ví dụ}, #1%
}

\begin{document}

\thispagestyle{empty}
\begin{tikzpicture}[remember picture, overlay]
  \fill[blue!80!black] (current page.south west) rectangle (current page.north east);
  \fill[blue!60!cyan, opacity=0.8]
    (current page.south west) -- (current page.north west) --
    ([xshift=-3cm]current page.north east) -- (current page.south east) -- cycle;
  \node[anchor=west, white, font=\fontsize{28}{34}\bfseries\selectfont]
    at ([xshift=3cm, yshift=-7.5cm]current page.north west) {GIẢI TÍCH 1};
  \node[anchor=west, white, font=\fontsize{20}{26}\selectfont]
    at ([xshift=3cm, yshift=-9cm]current page.north west) {Vi tích phân hàm một biến --- Năm 1, HK1};
  \node[anchor=west, white, font=\fontsize{18}{24}\selectfont]
    at ([xshift=3cm, yshift=-10cm]current page.north west) {Tài liệu lý thuyết};
  \node[anchor=west, white, font=\Large\bfseries]
    at ([xshift=3cm, yshift=-13cm]current page.north west) {Biên soạn:};
  \node[anchor=west, yellow!80!white, font=\fontsize{24}{30}\bfseries\selectfont]
    at ([xshift=3cm, yshift=-14.5cm]current page.north west) {ĐẶNG TẤN QUÍ};
  \node[anchor=south, white!80, font=\small]
    at ([yshift=1.5cm]current page.south) {Tài liệu có bản quyền --- Cấm sao chép khi chưa được sự đồng ý của tác giả};
  \node[anchor=south east, white, font=\Large\bfseries]
    at ([xshift=-2cm, yshift=3cm]current page.south east) {\the\year};
\end{tikzpicture}

\newpage
\tableofcontents
\newpage

%% ================================================================
%%  CHƯƠNG 1: HÀM SỐ, GIỚI HẠN, LIÊN TỤC
%% ================================================================
\section{Hàm số, giới hạn và liên tục}

\subsection{Hàm số}

\begin{dinhnghia}
  \textbf{Hàm số} $f: D \to \mathbb{R}$ là quy tắc cho tương ứng mỗi $x \in D$ với đúng một số thực $f(x)$. $D$ gọi là \textbf{tập xác định}.
\end{dinhnghia}

\begin{dinhnghia}
  Hàm số $f$ gọi là \textbf{chẵn} nếu $f(-x) = f(x)$, \textbf{lẻ} nếu $f(-x) = -f(x)$.
\end{dinhnghia}

\begin{dinhnghia}
  Hàm số $f$ gọi là \textbf{tuần hoàn} chu kỳ $T > 0$ nếu $f(x + T) = f(x)$ với mọi $x$.
\end{dinhnghia}

\subsection{Giới hạn dãy số}

\begin{dinhnghia}
  Dãy $(x_n)$ có \textbf{giới hạn} $L$ (viết $\lim_{n \to \infty} x_n = L$) nếu:
  \[
    \forall \varepsilon > 0,\; \exists N \in \mathbb{N}:\; n > N \Rightarrow |x_n - L| < \varepsilon.
  \]
\end{dinhnghia}

\begin{dinhly}[title={Giới hạn cơ bản}]
  \begin{itemize}
    \item $\displaystyle\lim_{n \to \infty} \frac{1}{n^p} = 0$ ($p > 0$).
    \item $\displaystyle\lim_{n \to \infty} q^n = 0$ nếu $|q| < 1$.
    \item $\displaystyle\lim_{n \to \infty} \left(1 + \frac{1}{n}\right)^n = e$.
  \end{itemize}
\end{dinhly}

\subsection{Giới hạn hàm số}

\begin{dinhnghia}
  $\displaystyle\lim_{x \to x_0} f(x) = L$ nếu:
  \[
    \forall \varepsilon > 0,\; \exists \delta > 0:\; 0 < |x - x_0| < \delta \Rightarrow |f(x) - L| < \varepsilon.
  \]
\end{dinhnghia}

\begin{congthuc}[title={Giới hạn quan trọng}]
  \begin{itemize}
    \item $\displaystyle\lim_{x \to 0} \frac{\sin x}{x} = 1$.
    \item $\displaystyle\lim_{x \to 0} \frac{e^x - 1}{x} = 1$.
    \item $\displaystyle\lim_{x \to \infty} \left(1 + \frac{1}{x}\right)^x = e$.
  \end{itemize}
\end{congthuc}

\subsection{Hàm số liên tục}

\begin{dinhnghia}
  Hàm số $f$ \textbf{liên tục} tại $x_0$ nếu $\displaystyle\lim_{x \to x_0} f(x) = f(x_0)$.
\end{dinhnghia}

\begin{dinhly}[title={Định lý giá trị trung gian}]
  Nếu $f$ liên tục trên $[a,b]$ và $f(a) \cdot f(b) < 0$ thì tồn tại $c \in (a,b)$ sao cho $f(c) = 0$.
\end{dinhly}

%% ================================================================
%%  CHƯƠNG 2: ĐẠO HÀM VÀ VI PHÂN
%% ================================================================
\section{Đạo hàm và vi phân}

\subsection{Đạo hàm}

\begin{dinhnghia}
  \textbf{Đạo hàm} của $f$ tại $x_0$:
  \[
    f'(x_0) = \lim_{h \to 0} \frac{f(x_0 + h) - f(x_0)}{h} = \lim_{x \to x_0} \frac{f(x) - f(x_0)}{x - x_0}.
  \]
\end{dinhnghia}

\begin{congthuc}[title={Đạo hàm cơ bản}]
  \begin{itemize}
    \item $(x^n)' = n x^{n-1}$.
    \item $(\sin x)' = \cos x$, $(\cos x)' = -\sin x$.
    \item $(e^x)' = e^x$, $(\ln x)' = \dfrac{1}{x}$.
    \item $(u \cdot v)' = u'v + uv'$.
    \item $\left(\dfrac{u}{v}\right)' = \dfrac{u'v - uv'}{v^2}$.
    \item $(f \circ g)'(x) = f'(g(x)) \cdot g'(x)$.
  \end{itemize}
\end{congthuc}

\subsection{Vi phân}

\begin{dinhnghia}
  \textbf{Vi phân} của $f$ tại $x$: $df = f'(x)\, dx$.
\end{dinhnghia}

\begin{dinhly}[title={Định lý Fermat}]
  Nếu $f$ đạt cực trị địa phương tại $x_0$ và $f'(x_0)$ tồn tại thì $f'(x_0) = 0$.
\end{dinhly}

\begin{dinhly}[title={Định lý Lagrange}]
  Nếu $f$ liên tục trên $[a,b]$ và có đạo hàm trên $(a,b)$ thì tồn tại $c \in (a,b)$:
  \[
    f'(c) = \frac{f(b) - f(a)}{b - a}.
  \]
\end{dinhly}

\subsection{Khai triển Taylor}

\begin{congthuc}
  \textbf{Khai triển Taylor} của $f$ tại $x_0$:
  \[
    f(x) = f(x_0) + f'(x_0)(x - x_0) + \frac{f''(x_0)}{2!}(x - x_0)^2 + \cdots + \frac{f^{(n)}(x_0)}{n!}(x - x_0)^n + R_n(x).
  \]
  Khi $x_0 = 0$ gọi là \textbf{khai triển Maclaurin}.
\end{congthuc}

\begin{vidu}
  $e^x = 1 + x + \dfrac{x^2}{2!} + \dfrac{x^3}{3!} + \cdots$, $\quad \sin x = x - \dfrac{x^3}{3!} + \dfrac{x^5}{5!} - \cdots$.
\end{vidu}

%% ================================================================
%%  CHƯƠNG 3: TÍCH PHÂN
%% ================================================================
\section{Tích phân}

\subsection{Nguyên hàm}

\begin{dinhnghia}
  $F$ gọi là \textbf{nguyên hàm} của $f$ trên $I$ nếu $F'(x) = f(x)$ với mọi $x \in I$. Ký hiệu $\displaystyle\int f(x)\,dx = F(x) + C$.
\end{dinhnghia}

\begin{congthuc}[title={Nguyên hàm cơ bản}]
  \begin{itemize}
    \item $\displaystyle\int x^n\,dx = \frac{x^{n+1}}{n+1} + C$ ($n \ne -1$).
    \item $\displaystyle\int \frac{1}{x}\,dx = \ln|x| + C$.
    \item $\displaystyle\int e^x\,dx = e^x + C$.
    \item $\displaystyle\int \sin x\,dx = -\cos x + C$, $\displaystyle\int \cos x\,dx = \sin x + C$.
  \end{itemize}
\end{congthuc}

\subsection{Tích phân xác định}

\begin{dinhnghia}
  \textbf{Tích phân xác định}: $\displaystyle\int_a^b f(x)\,dx = F(b) - F(a)$, với $F$ là nguyên hàm của $f$.
\end{dinhnghia}

\begin{dinhly}[title={Định lý cơ bản của giải tích}]
  Nếu $f$ liên tục trên $[a,b]$ và $F(x) = \displaystyle\int_a^x f(t)\,dt$ thì $F'(x) = f(x)$.
\end{dinhly}

\begin{phuongphap}[title={Phương pháp đổi biến}]
  Đặt $x = \varphi(t)$ với $\varphi$ khả vi, đơn điệu. Khi đó $\displaystyle\int f(x)\,dx = \int f(\varphi(t))\,\varphi'(t)\,dt$.
\end{phuongphap}

\begin{phuongphap}[title={Tích phân từng phần}]
  $\displaystyle\int u\,dv = uv - \int v\,du$.
\end{phuongphap}

%% ================================================================
%%  CHƯƠNG 4: CHUỖI SỐ
%% ================================================================
\section{Chuỗi số}

\begin{dinhnghia}
  \textbf{Chuỗi số} $\displaystyle\sum_{n=1}^{\infty} a_n$: tổng vô hạn $a_1 + a_2 + \cdots$. Tổng riêng $S_n = a_1 + \cdots + a_n$. Chuỗi \textbf{hội tụ} nếu $\lim_{n \to \infty} S_n$ tồn tại hữu hạn.
\end{dinhnghia}

\begin{dinhly}[title={Điều kiện cần hội tụ}]
  Nếu $\displaystyle\sum a_n$ hội tụ thì $\lim_{n \to \infty} a_n = 0$. (Điều ngược lại không đúng.)
\end{dinhly}

\begin{dinhly}[title={Chuỗi dương}]
  \begin{itemize}
    \item \textbf{Tiêu chuẩn so sánh}: Nếu $0 \le a_n \le b_n$ và $\sum b_n$ hội tụ thì $\sum a_n$ hội tụ.
    \item \textbf{Tiêu chuẩn D'Alembert}: $\displaystyle\lim_{n \to \infty} \frac{a_{n+1}}{a_n} = L$: $L < 1$ hội tụ, $L > 1$ phân kỳ.
    \item \textbf{Tiêu chuẩn Cauchy}: $\displaystyle\lim_{n \to \infty} \sqrt[n]{a_n} = L$: $L < 1$ hội tụ, $L > 1$ phân kỳ.
  \end{itemize}
\end{dinhly}

\begin{congthuc}
  Chuỗi hình học: $\displaystyle\sum_{n=0}^{\infty} q^n = \frac{1}{1-q}$ khi $|q| < 1$.
\end{congthuc}

\end{document}
